\section{方法的优势与局限}

\subsection{推理时引导方法的全面比较}

\begin{center}
\begin{tabular}{p{3cm}p{3.5cm}p{3.5cm}p{3.5cm}}
\toprule
\textbf{方法} & \textbf{注意力操纵} & \textbf{DPS (逆问题)} & \textbf{奖励函数引导} \\
\midrule
理论基础 & 无（工程直觉） & 贝叶斯法则 & RL / KL 约束优化 \\
适用条件 & 特定架构 & 已知映射 $\mathcal{A}$ & 可微奖励函数 \\
通用性 & 低 & 中 & 高 \\
实现难度 & 中 & 低 & 低 \\
训练需求 & 无 & 无 & 无（或少量） \\
典型应用 & 布局控制 & 超分/去噪/修复 & 任意偏好引导 \\
\bottomrule
\end{tabular}
\end{center}

\subsection{优势}

\begin{enumerate}
    \item \textbf{通用性}：只要能定义可微的奖励函数，就可以应用。适用于图像、视频、运动、分子等多种模态。
    \item \textbf{无需额外训练}：利用现有预训练模型，无需收集新数据或进行额外训练。
    \item \textbf{实现简单}：核心代码只需在生成流水线中添加梯度引导步骤。
    \item \textbf{灵活组合}：可以同时施加多个奖励函数，实现多条件引导。
    \item \textbf{计算友好}：相比从头训练，推理时引导只需较少的计算资源。
\end{enumerate}

\subsection{局限}

\begin{warningbox}{奖励函数必须可微}
推理时引导依赖于梯度反传，因此奖励函数必须是可微的。对于不可微的评价指标（如 FID、IS 等），不能直接用作引导信号。这是当前方法的一个根本限制。
\end{warningbox}

\begin{warningbox}{真实性与条件满足的权衡不可避免}
增大引导强度会使中间样本偏离训练分布，导致噪声预测网络在分布外区域产生不准确的预测。最终表现为：
\begin{itemize}
    \item 引导强度小：输出真实但不满足条件
    \item 引导强度大：满足条件但输出不真实
\end{itemize}
寻找最佳引导权重通常需要经验性调参。
\end{warningbox}

其他局限包括：
\begin{itemize}
    \item 推理速度变慢：更多时间步能提供更好的引导效果，但增加了推理时间
    \item Tweedie 近似在高噪声步误差较大：早期时间步的 $\hat{x}_0$ 估计质量差
    \item 没有理论保证收敛到全局最优
\end{itemize}

\begin{knowledgebox}{与训练时方法的对比}
推理时引导与训练时方法（如 ControlNet、条件扩散模型、RLHF 微调）形成互补：
\begin{itemize}
    \item \textbf{训练时方法}：效果更稳定、生成速度快，但需要大量数据和计算资源进行训练，且每种条件需要单独训练
    \item \textbf{推理时方法}：无需训练、灵活性高、可即时组合多种条件，但效果不够稳定、生成速度较慢
\end{itemize}
在实际工程中，最佳策略往往是二者结合：用训练时方法处理常见的、固定的条件类型，用推理时引导处理临时的、多样化的条件需求。
\end{knowledgebox}

\subsection{本章小结}

推理时引导方法具有通用性强、无需训练、实现简单等优势，但受限于奖励函数必须可微以及真实性-条件满足度的固有权衡。

